% =====================================================================
% VERSION ENSAE, REDUCTION SFE (phase 4, 1er octobre 2026) : notes de synthese.
% Consignes MS 2025-2026 du service des stages : deux notes autonomes, francaise puis
% anglaise, d'une a deux pages chacune, en fin de fichier. La version de quatre pages
% chacune, avec la figure LUCY et la figure N1, reste celle de la version Institut
% (chapitres/20_notes_synthese.tex) et l'historique Git de ce fichier.
% Lu par main_ensae.tex seul : ne jamais le faire lire par main_institut.tex.
%
% Chaque nombre est une valeur publiee dans le corps ; la note n'en calcule aucun. Elle cite
% ses scripts comme une section, pour rester sous harnais.
% =====================================================================

\providecommand*{\notesec}[1]{\par\smallskip\noindent{\bfseries #1}\par}

\chapter*{Note de synthèse}
\addcontentsline{toc}{chapter}{Note de synthèse}
\markboth{Note de synthèse}{}

% SOURCES-SCRIPTS: 28 30 43 48 60 63 67 68 74 91 92 95 110 08b 89 99 65 66 76 38

\notesec{Contexte et objectif}
Depuis janvier 2025, le règlement européen DORA impose aux banques et aux assureurs de maîtriser
cinq domaines de résilience numérique, que nous appelons piliers : la gouvernance du risque
informatique (P1), la gestion des incidents (P2), les tests (P3), les prestataires tiers (P4) et le
partage d'informations (P5). Le règlement et Solvabilité~II ne disent pas ce que coûte une non-conformité : la charge de risque opérationnel de la Formule Standard est la même, qu'une entité soit conforme ou non. Dans ce stage, mené chez Nexialog
Consulting sous la direction de Hugo Rapior, nous chiffrons en capital l'écart entre ces deux
situations.

Le risque mesuré est le risque opérationnel d'origine informatique de l'entité. Les pénalités réglementaires et le risque de souscription des polices cyber
qu'un assureur vend ne sont pas modélisés. Nous répartissons l'écart entre les cinq piliers et entre
quatre paramètres du modèle que la non-conformité modifie, que nous appelons canaux : la fréquence
des incidents, leur détection, leur propagation d'un pilier à l'autre et la dépendance à un
prestataire commun. Le résultat est un besoin de capital au sens de l'ORSA (évaluation interne des risques), non un SCR réglementaire.

Les montants sont en millions d'euros et valent pour le périmètre financier étudié ($21{,}6$ incidents par an), non pour un assureur ; l'ordre de grandeur d'une entité de $15\,000$~M\euro{} de provisions figure avec les résultats.
Le niveau change d'un facteur $22$ selon la seule source de sévérité : nous interprétons l'écart
entre les deux états et l'ordre des piliers, pas le niveau.

\notesec{Démarche}
La base SAS OpRisk Global fournit $582$ pertes informatiques d'institutions financières. La fréquence conserve un dénominateur conventionnel de $27$ ans, à réconcilier avec les événements antérieurs à 2000. Une loi de Pareto généralisée, ajustée au-delà de $20{,}03$~M\euro{} sur $91$ excès, décrit les
pertes élevées (indice de queue $0{,}5954$). Le nombre annuel d'incidents suit une loi binomiale négative. Un incident naît sur
un pilier, puis se transmet aux autres selon une matrice d'influence dirigée, posée à dire d'expert
puis bornée.  Le
capital est le quantile à $99{,}5\,\%$ de la charge annuelle, calculé dans chaque état.

La non-conformité n'ajoute pas de pénalité : elle change la valeur de quatre paramètres. La fréquence est multipliée par $2{,}49$ : les parts des vecteurs d'attaque de la base
Hackmageddon sont pondérées par des surcroîts d'incidents tirés de résultats publiés.
La détection agit sur chaque pilier touché. Nous multiplions sa probabilité de sévérité supérieure à
$20{,}03$~M\euro{} par $0{,}85$ à l'état conforme et par $1{,}20$ à l'état non conforme.
La probabilité obtenue est de $0{,}128$ contre $0{,}181$ par pilier touché.

Le gain de la cascade, nombre de piliers
que le pilier le plus contagieux touche directement, vaut $0{,}45$ à l'état conforme et $0{,}90$ à
l'état non conforme. L'accumulation chez un prestataire commun passe de $0$ à
$0{,}68$, valeur fixée d'après un indicateur de concentration des fournisseurs de cloud.

Ces écarts entre états sont posés. Cinq correspondances état--gain donnent le même classement. La fréquence est testée de $41{,}7$ à $65{,}6$ incidents/an. L'accumulation fait varier le capital de $7\,694$ à $10\,096$~M\euro{} entre $0$ et $0{,}68$. Les multiplicateurs de détection ne sont pas balayés ; seule l'élasticité au taux de dépassement, $0{,}82$, est mesurée.

\begin{center}\footnotesize
\renewcommand{\arraystretch}{1.1}
\begin{tabular}{@{}l l r r r r@{}}
\toprule
\textbf{Canal} & \textbf{Paramètre} & \textbf{Conforme} & \textbf{Non conforme} &
\textbf{Effet isolé} & \textbf{Effet de fermeture} \\
\midrule
Fréquence & incidents par an & $21{,}6$ & $53{,}6$ & $4\,328$ & $8\,775$ \\
Détection & dépassement par pilier & $0{,}128$ & $0{,}181$ & $1\,448$ & $4\,562$ \\
Propagation & gain de la cascade & $0{,}45$ & $0{,}90$ & $1\,633$ & $2\,402$ \\
Accumulation & prestataire commun & $0$ & $0{,}68$ & $1\,728$ & $3\,402$ \\
\midrule
\multicolumn{2}{@{}l}{\textbf{Besoin de capital} (M\euro)} & $6\,049$ & $20\,188$ & & \\
\bottomrule
\end{tabular}
\end{center}

\notesec{Résultats}
À l'échelle du périmètre financier étudié, le besoin de capital est multiplié par $3{,}34$ entre
l'entité conforme et l'entité non conforme sur les quatre canaux, soit un écart de $14\,139$~M\euro.
Les canaux pris un à un en expliquent $9\,138$~M\euro, et leurs interactions $5\,001$~M\euro, soit
$35\,\%$. L'effet isolé est le capital ajouté par un canal seul, depuis l'état conforme. L'effet de
fermeture est le capital retiré en corrigeant ce canal seul, depuis l'état non conforme sur les trois
autres : un plan de remédiation doit citer ce second effet.

La fréquence seule (le nombre d'incidents informatiques) porte le capital
à $10\,377$~M\euro, alors que les trois autres canaux réunis n'atteignent que $11\,413$~M\euro. Les données publiques ne disent pas quel pilier entraîne l'autre : nous calculons le capital aux
$1\,024$ sommets de l'ensemble des matrices d'influence admissibles, la structure symétrique étant
posée à dire d'expert. Dans le calage de référence (fréquence conforme, gain $0{,}90$), les valeurs aux sommets vont de $6\,858$ à
$8\,697$~M\euro.  Pour une entité notionnelle de $15\,000$~M\euro{} de provisions, le capital du secteur est
ramené à sa taille : le besoin est de l'ordre de $169$~M\euro, une extrapolation d'ordre de grandeur. Ces valeurs dépendent de la structure posée à dire d'expert et de la calibration.

\notesec{Limites et recommandations}
La propagation et l'accumulation sont posées. Pour la propagation, les correspondances testées conservent le signe de l'écart ; son amplitude dépend de ces choix. L'hypothèse la plus lourde est que les coûts des
piliers touchés par un même incident s'additionnent.

Le choix de la loi binomiale négative s'appuie sur un test du rapport de vraisemblance, qui rejette
la loi de Poisson (statistique de $130{,}5$). Un backtest sur $2004$--$2025$ place ensuite $11$
années sur $12$ dans l'intervalle à $90\,\%$ de la loi binomiale négative, contre $9$ sur $12$ pour
Poisson, ce qui conforte le choix sans le démontrer sur douze années. Le backtest détecte aussi une dérive de l'échelle des pertes, qui rend le niveau publié plutôt trop bas. Le quantile à $99{,}5\,\%$ ne peut être testé sur aucun historique existant : il faudrait $1\,811$ années.

Pour l'entité, la priorité est de maîtriser la fréquence de ses incidents informatiques, et de chiffrer toute remédiation par son effet de fermeture. Pour le
régulateur, un registre d'incidents daté, classé par domaine de contrôle et consignant tous les
piliers touchés rendrait la direction de la propagation mesurable. VERIS prévoit déjà ces champs ; ils sont renseignés ensemble dans un incident sur $10\,591$.

% =====================================================================

\chapter*{Executive summary}
\addcontentsline{toc}{chapter}{Executive summary}
\markboth{Executive summary}{}

% SOURCES-SCRIPTS: 28 30 43 48 60 63 67 68 74 91 92 95 110 08b 89 99 65 66 76 38

\begin{otherlanguage*}{english}
\notesec{Context and objective}
Since January 2025, the European DORA regulation has required banks and insurers to master five
areas of digital operational resilience, which we call pillars: IT risk governance (P1), incident
management (P2), testing (P3), third-party providers (P4) and information sharing (P5). The
regulation does not say what non-compliance costs, and neither does Solvency~II: the Standard
Formula operational risk charge is the same for a compliant entity and for one that is not. In this
internship, carried out at Nexialog Consulting under the supervision of Hugo Rapior, we express in
capital terms the gap between these two situations.

The risk measured is the entity's own IT-related operational risk, that is its own incidents.
Regulatory penalties and the underwriting risk of the cyber policies an insurer sells are not
modelled. We split the gap across the five pillars and across four model parameters that
non-compliance changes, which we call channels: incident frequency, detection, propagation from one
pillar to another, and dependence on a shared provider. The result is a capital need in the ORSA
sense, the internal assessment of risks and solvency. It is not a regulatory SCR.

The amounts are in millions of euros and are for the financial sample taken as a whole
($21.6$ incidents a year in the loss database), not for one insurer; the order of magnitude for an
entity with $15\,000$~M\euro{} of provisions is given with the results. The level changes by a factor of $22$ with the severity source alone: we interpret the gap
between the two states and the ordering of pillars, not the level.

\notesec{Approach}
SAS OpRisk Global provides $582$ IT losses at financial institutions. Frequency retains a conventional $27$-year denominator, which requires reconciliation with events before 2000. A generalised Pareto distribution, fitted above $20.03$~M\euro{} on $91$ exceedances, describes
the high losses (tail index $0.5954$). The annual
incident count is negative binomial. An incident starts in one pillar, then spreads to the others
through a directed influence matrix, posited by expert judgement and then bounded.  Capital is the $99.5\,\%$ quantile
of the annual loss, computed in each state.

Non-compliance does not add a penalty: it changes the value of four parameters. Frequency is multiplied by $2.49$: the attack-vector shares of the Hackmageddon database are
weighted by excess incident rates taken from published results. Detection acts on each affected pillar.
We multiply its probability of a severity above $20.03$~M\euro{} by $0.85$ in the compliant state and by $1.20$ in the non-compliant state.
This gives $0.128$ against $0.181$ per affected pillar.

The cascade gain, the number of pillars directly hit by the most contagious
pillar, is $0.45$ in the compliant state and $0.90$ in the non-compliant one.
Accumulation at a shared provider goes from $0$ to $0.68$, a value set from a cloud-provider
concentration indicator.

These differences between states are posited. Five state--gain mappings preserve the ranking. Frequency is tested from $41.7$ to $65.6$ incidents/year. Accumulation takes capital from $7\,694$ to $10\,096$~M\euro{} between $0$ and $0.68$. Detection multipliers are not swept; only the elasticity to the exceedance rate, $0.82$, is measured.

\begin{center}\footnotesize
\renewcommand{\arraystretch}{1.1}
\begin{tabular}{@{}l l r r r r@{}}
\toprule
\textbf{Channel} & \textbf{Parameter} & \textbf{Compliant} & \textbf{Non-compliant} &
\textbf{Standalone} & \textbf{Closure} \\
\midrule
Frequency & incidents per year & $21.6$ & $53.6$ & $4\,328$ & $8\,775$ \\
Detection & exceedance per pillar & $0.128$ & $0.181$ & $1\,448$ & $4\,562$ \\
Propagation & cascade gain & $0.45$ & $0.90$ & $1\,633$ & $2\,402$ \\
Accumulation & shared provider & $0$ & $0.68$ & $1\,728$ & $3\,402$ \\
\midrule
\multicolumn{2}{@{}l}{\textbf{Capital need} (M\euro)} & $6\,049$ & $20\,188$ & & \\
\bottomrule
\end{tabular}
\end{center}

\notesec{Results}
At the scale of the financial sample, capital need is multiplied by $3.34$ between the
compliant entity and the entity that is non-compliant on all four channels, a gap of $14\,139$~M\euro.
The channels taken one by one explain $9\,138$~M\euro, and their interactions $5\,001$~M\euro, or
$35\,\%$. The standalone effect is the capital added by one channel alone, from the compliant state. The
closure effect is the capital removed by fixing that channel alone, from the state that is
non-compliant on the other three: a remediation plan should quote the second.

Frequency alone, that is the number of IT incidents at the entity, takes capital to
$10\,377$~M\euro, while the three other channels together reach only $11\,413$~M\euro. Public data do not say which pillar drives the other: we compute capital at the $1\,024$
vertices of the set of admissible influence matrices, the symmetric structure being set by expert
judgement. In the reference calibration (compliant-state frequency, gain $0.90$),
the vertex values range from $6\,858$ to $8\,697$~M\euro.  For a notional entity with $15\,000$~M\euro{} of provisions, sector capital is scaled down to its
size: the need is about $169$~M\euro, an extrapolated order of magnitude. These values depend on the structure set by expert judgement and on the calibration.

\notesec{Limitations and recommendations}
Propagation and accumulation are posited. For propagation, the tested mappings preserve the sign of the gap; its size depends on these choices. The heaviest assumption is that the costs of the pillars hit
by one incident add up.

The choice of the negative binomial rests on a likelihood ratio test that rejects the Poisson
distribution (statistic of $130.5$). A backtest over $2004$--$2025$ then places $11$ years out of
$12$ inside the $90\,\%$ interval of the negative binomial, against $9$ out of $12$ for the Poisson.
This supports the choice without proving it over twelve years. The backtest also detects a drift in loss scale, which makes the published level probably too low. The $99.5\,\%$ quantile cannot be tested on any existing history: it would take $1\,811$ years.

For the entity, the priority is to control the number of IT incidents, that is the frequency of its
operational risk, and to cost any remediation by its closure effect. For the regulator, a dated
incident register, classified by control domain and recording every pillar hit, would make the
direction of propagation measurable. VERIS already provides these fields; they are filled in together in only one incident out of
$10\,591$.
\end{otherlanguage*}
